\section{Methods}
\label{sec:methods}

\subsection{Pipeline Overview}
\label{subsec:pipeline_overview}

The measures below come from one integrated pipeline, not independent
scripts. Lifting the cylinder off its center peg triggers the Arduino
board signal (Section~\ref{subsec:task}); this starts the GUI trial
clock and begins path-length accumulation. Placing the cylinder on its
target peg triggers the matching Arduino signal and stops the clock. At
that instant, HSV segmentation on the YOLOv8-detected~\cite{yolov8} cylinder estimates
its orientation and checks whether the correct face -- blue or white, as
indicated by the lit target peg (Section~\ref{subsec:task}) -- ended
face down. Once the trial ends, the GUI plots the total force applied to
the task pad during the trial: time on the x-axis, force in newtons on
the y-axis, read from the ATI sensor.

Table~\ref{tab:data_timing} summarizes every measure the pipeline
collects: its timing class, what triggers it, and how it reaches the
trainee.

\begin{table}[h]
\centering
\resizebox{\textwidth}{!}{%
\begin{tabular}{|p{5.2cm}|p{3.5cm}|p{5.6cm}|p{5.6cm}|}
\hline
\textbf{Data / Measure} & \textbf{Timing} & \textbf{Trigger} &
\textbf{Feedback} \\
\hline
Trial time &
Real-time &
Continuous, lift-to-placement &
Progress bar (GUI) \\[4pt]
\hline
Path length &
Real-time &
Continuous while clock runs &
Numeric running total \\[4pt]
\hline
Object detection (YOLOv8~\cite{yolov8}) &
Real-time, delayed &
Continuous, event-triggered at placement &
Backbone for orientation \& state below \\[4pt]
\hline
Cylinder orientation (HSV) &
Real-time, delayed &
Continuous, event-triggered at placement &
Binary indicator \\[4pt]
\hline
Object state (4 classes) &
Real-time &
Continuous &
Status indicator \\[4pt]
\hline
Contact force &
Real-time, delayed &
Continuous, after trial ends &
2D graph, time (x) vs.\ force in N (y) \\[4pt]
\hline
Additional camera &
Real-time &
Continuous, pedal-triggered &
Additional camera view \\[4pt]
\hline
3D localization (SAM2~\cite{sam2}) &
Post-hoc &
Offline processing &
Not shown to trainee; offline plot \\
\hline
\end{tabular}%
}
\caption{Data collected by the pipeline, its timing class, trigger, and
the feedback form it reaches the trainee in}
\label{tab:data_timing}
\end{table}

\subsection{Real-Time GUI \& Trial Clock}

The trial clock starts and stops from the Arduino cylinder-lift/placement
signal (Section~\ref{subsec:pipeline_overview}), not from pedal state.
Path length accrues on the GUI for the full duration the clock runs.
% TODO: code to be copied into code/ from external folder once path is
% provided.

\subsection{Object Detection}
% TODO: YOLOv8 detector, trained on 4 classes -- cylinders, inactive
% pegs, blue pegs, white pegs. Feeds the angular-displacement and
% object-state methods below.

\subsection{Path Length}
% TODO

\subsection{Angular Displacement}
\label{subsec:angular_displacement}

Cylinder orientation is estimated by HSV segmentation on the region
YOLOv8 detects as the cylinder.
% TODO: describe the HSV method itself -- color thresholds, how the
% orientation angle is derived from the segmented mask.

This estimate is computed at placement and feeds the
orientation-correctness check in Object State Detection below.

\subsection{Object State Detection}

At placement, the orientation estimate from
Section~\ref{subsec:angular_displacement} determines whether the
cylinder landed correctly: the lit target peg (Section~\ref{subsec:task})
specifies blue or white face down, and the trial is scored correct or
incorrect against that target.
% TODO: classifier over 4 cylinder states -- stationary, held, dropped,
% lost -- augmented with XGBoost. Cross-check against conductive-paint
% task board signal (subsec:task) and ground contact detection below.

\subsection{Ground Contact Detection}

Contact force from the ATI sensor is not streamed live during the trial.
Once the trial ends, the GUI displays it as a 2D graph -- time on the
x-axis, total force in newtons on the y-axis -- showing how much force
was applied to the task pad over the course of the trial.

\subsection{3D Localization (Post-Hoc)}

The pipeline includes a 3D localization method for the cylinder, tracking
its $(x, y, z)$ position from the RealSense D405 depth stream. An early
version of this method ran in real time, but tracking was unreliable.
% TODO: state the specific failure mode (noise / occlusion / drift / other)
To improve accuracy, the method was rebuilt on Segment Anything 2
(SAM2)~\cite{sam2} for cylinder segmentation. SAM2 produces an accurate
mask, but its per-frame runtime exceeds the trial's real-time budget.

This method therefore runs post-hoc, after the trial ends, not live
during it. Its role in this thesis is validation, not live measurement:
3D localization gives an accurate offline reference for cylinder
position, used to check the real-time measures -- path length and
angular displacement -- against ground truth. It is not one of the
real-time measures the system delivers during a trial
(Section~\ref{subsec:scope}).